\documentclass[11pt,a4paper]{article}
\usepackage[utf8]{inputenc}
\usepackage[T1]{fontenc}
\usepackage[italian]{babel}
\usepackage{amsmath,amssymb}
\usepackage{geometry}
\usepackage{hyperref}
\geometry{margin=2.7cm}
\title{Ipotesi della Curvatura Minima Universale I:\\fondamenti cinematici e criteri di confutazione}
\author{Pepe Donato\\Independent Researcher\\\texttt{donato.pepe.it@gmail.com}}
\date{Bozza auditata --- 2026}

\begin{document}
\maketitle

\begin{abstract}
Studiamo come ipotesi falsificabile l'esistenza di una scala universale $\kappa_0>0$ con dimensione inversa di lunghezza. Questa bozza non presenta una rivelazione né una teoria completa. L'audit del documento sorgente sostiene, con condizioni, la cinematica di Frenet--Serret per curve timelike, ma non deriva il limite positivo. Con la definizione $\kappa=\sqrt{a^\mu a_\mu}/c^2$, una geodetica timelike ha $\kappa=0$; quindi il postulato è in tensione con la caduta libera ideale finché non viene fornita una nuova dinamica covariante compatibile con esperimenti. Stato del postulato: \texttt{UNPROVEN}.
\end{abstract}

\section{Stato epistemico}\label{sec:status}
Il progetto separa risultati standard, definizioni, ipotesi e conseguenze. I claim sorgente UMCH-CLM-0008, UMCH-CLM-0009, UMCH-CLM-0010 e UMCH-CLM-0011 sono stati auditati insieme ai duplicati UMCH-CLM-0049, UMCH-CLM-0050, UMCH-CLM-0051, UMCH-CLM-0052 e alla riformulazione UMCH-CLM-0092, UMCH-CLM-0093, UMCH-CLM-0094. Il formalismo per curve timelike è documentato in \cite{FormigaRomero2006}; il quadro sperimentale relativistico è discusso in \cite{Will2014}. Nessuna fonte citata dimostra $\kappa_0>0$.

\section{Definizioni e dominio}\label{sec:definitions}
Sia $x^\mu(\tau)$ una curva timelike regolare in una varietà lorentziana con segnatura $(-,+,+,+)$. Poniamo
\begin{equation}\label{eq:velocity}
u^\mu=\frac{dx^\mu}{d\tau},\qquad u^\mu u_\mu=-c^2,
\end{equation}
e
\begin{equation}\label{eq:acceleration}
a^\mu=u^\nu\nabla_\nu u^\mu.
\end{equation}
Quando $a^\mu\ne0$, definiamo curvatura propria della linea d'universo e normale
\begin{equation}\label{eq:curvature}
\kappa=\frac{\sqrt{a^\mu a_\mu}}{c^2},\qquad N^\mu=\frac{a^\mu}{c^2\kappa}.
\end{equation}
Non identifichiamo questa quantità con curvatura intrinseca dello spaziotempo. Curve null, campi e vuoto non sono coperti da questa definizione.

\section{Ipotesi UMCH}\label{sec:hypothesis}
Il nucleo congettura
\begin{equation}\label{eq:hypothesis}
\kappa\ge\kappa_0>0,\qquad [\kappa_0]=L^{-1}.
\end{equation}
$\kappa_0$ resta parametro libero. Relazioni con $c,\hbar,G,\Lambda$ o scale cosmologiche non appartengono al nucleo senza derivazione indipendente. UMCH-CLM-0093 resta \texttt{UNPROVEN}.

\section{Derivazione cinematica minima}\label{sec:derivation}
Per compatibilità metrica,
\begin{equation}\label{eq:orthogonality}
\frac{D}{d\tau}(u^\mu u_\mu)=2u_\mu a^\mu=0.
\end{equation}
Se $\kappa\ne0$, le definizioni danno
\begin{equation}\label{eq:frenet}
a^\mu=c^2\kappa N^\mu,\qquad N^\mu N_\mu=1.
\end{equation}
Questa è identità definitoria, non prova della disuguaglianza \eqref{eq:hypothesis}. Per $P^\mu=m_0u^\mu$ con massa costante,
\begin{equation}\label{eq:momentum}
P^\mu P_\mu=-m_0^2c^2.
\end{equation}
L'ortogonalità conserva la norma locale di $u^\mu$; non implica da sola conservazione globale dell'energia-impulso.

\section{Tensione con moto geodetico}\label{sec:tension}
Una geodetica timelike parametrizzata dal tempo proprio soddisfa $a^\mu=0$. Da \eqref{eq:curvature} segue $\kappa=0$, e $N^\mu$ non è determinato. Quindi \eqref{eq:hypothesis} esclude la geodetica ideale. Una teoria UMCH deve introdurre un'azione o equazioni nuove, mostrare il limite standard e confrontarsi con test di equivalenza e Lorentz \cite{Will2014}. Curvatura non nulla dello spaziotempo non risolve questa tensione: spaziotempi curvi ammettono geodetiche.

\section{Limiti richiesti}\label{sec:limits}
Nel limite $\kappa_0\to0$ una deformazione proposta deve recuperare moto inerziale in relatività speciale, caduta geodetica in relatività generale e relativi principi di conservazione. Limiti null, deboli e non relativistici richiedono trattamenti espliciti. Un limite singolare deve essere derivato, non dichiarato.

\section{Criteri di confutazione}\label{sec:falsification}
La formulazione timelike viene respinta o ristretta se: non esiste dinamica covariante; compaiono instabilità o acausalità inevitabili; ogni $\kappa_0>0$ è escluso; il limite standard fallisce; oppure il modello non produce contenuto osservabile distinto. Prima di una misura positiva si cercano limiti superiori. Senza mappa derivata tra osservabile sperimentale e $\kappa_0$, non è lecito riportare un bound numerico.

\section{Limitazioni}\label{sec:limitations}
Questo lavoro non definisce ancora settori null, campi o vuoto; non propone un'azione; non dimostra stabilità; non tratta ALD, QED, gravità quantistica o cosmologia; non misura $\kappa_0$. La bibliografia è iniziale. Claim non auditati restano fuori dai risultati.

\section{Assistenza IA e responsabilità}\label{sec:ai}
Strumenti IA hanno assistito estrazione, classificazione, traduzione e redazione. Non sono fonti né coautori. Metadati bibliografici, formule, test e conclusioni richiedono verifica umana. Responsabilità scientifica resta all'autore.

\bibliographystyle{plain}
\bibliography{../../../references/library}
\end{document}
